\begin{topic}[id=printing_3d,
             difficulty={Einstieg},
             type={Überblick + Anwendungsanalyse},
             prerequisites={Allgemeines Technikinteresse; keine speziellen Vorkenntnisse erforderlich},
             practice={optional},
             owner={Dozententeam},
             updated={2026-04-09},
             tags={ Additive Fertigung, FDM, SLS, SLA, Material, Topologieoptimierung }]{3D-Druck: Additive Fertigung der Zukunft}

\TopicDescription{Additive Verfahren erlauben komplexe Geometrien, leichte Bauteile und schnelle Iteration. Wir vergleichen zentrale Technologien (FDM, SLS, SLA), Materialeigenschaften, Nachbearbeitung und typische Fehlerbilder. Dazu kommen Wirtschaftlichkeit, Nachhaltigkeitsaspekte und Anwendungsfelder in Industrie/Medizin. Ziel: pragmatische Entscheidungsgrundlagen für Technologie- und Materialwahl.}

\TopicLeitfragen{\item Welche Technologie passt zu welchem Bauteil?
\item Wie beeinflussen Material/Nachbearbeitung die Eigenschaften?
\item Wann lohnt additive Fertigung wirtschaftlich?
}
\TopicScope{\item Kein CAD/Designkurs.
\item Keine Drucker-spezifischen Hacks.
}
\TopicPractice{Optionale Praxisidee: Vergleich zweier Druckverfahren für dasselbe Bauteil (Eigenschaften/Kosten).}

\TopicKeywords{Additive Fertigung, FDM, SLS, SLA, Material, Topologieoptimierung}

\nocite{astmF42,3dPrintGuide}
\end{topic}
